% ── Chapter 2: Literature Review ─────────────────────────────

\chapter{Literature Review}

\section{Introduction}

This chapter traces the development of ideas and systems that are
directly relevant to the framework built in this work. The review
covers four interrelated areas: vision-based mobile robot navigation,
communication frameworks for distributed systems, control strategies,
and IoT-based environmental monitoring. The goal is not just to
catalogue what others have done but to understand where those efforts
stopped short and why the proposed system needed to be built the way
it was.

The literature is organised roughly chronologically, which helps show
how each area matured over time and why the gaps that remained in 2024
were not incidental but structural, being the result of researchers working
within separate disciplines that had limited cross-pollination.

\section{Early Developments in Vision-Based Navigation and Control}

The earliest serious work on vision-guided navigation for mobile robots
appeared in the late 1990s, when image-based visual servoing and
nonlinear control were still genuinely new ideas for this application.
Ma et al.\ \cite{r5} laid important theoretical groundwork with a
navigation system for nonholonomic mobile robots that used image-plane
features to close the control loop. The limitation was that the work
stayed in simulation; there was no physical robot, and the gap between
simulation and real hardware was left unaddressed.

Fang et al.\ \cite{r6} took the homography-based approach further,
developing stable control laws using Lyapunov methods. The mathematics
were rigorous, but again, experimental validation was absent. A later
paper by the same group \cite{r7} introduced adaptive control
strategies and included physical experiments, which was a significant
step forward, although the experiments required structured environments
with pre-placed visual markers, not something one would find in a
general indoor space.

Chen et al.\ \cite{r8} contributed a homography-based tracking control
method that improved trajectory following under disturbances. Jean and
Lian \cite{r14} later improved robustness by using geometric shape
parameters rather than point features, which reduced sensitivity to
partial occlusions. Both contributions, though valuable, were tied to
controlled settings and predefined feature structures.

Looking at these early works together, the pattern is clear: the focus
was almost entirely on the navigation and control problem. Environmental
sensing was not part of the picture, and communication was handled
implicitly through wired connections in lab setups.

\section{Communication Frameworks for Distributed Systems}

Papadimitratos et al.\ \cite{r3} proposed a secure vehicular
communication architecture that addressed reliability and security in
dynamic, multi-node environments. This was largely a networking
contribution; the sensing and application layers were outside its
scope, but the architecture principles carry over directly to
distributed IoT deployments.

Research on LTE and IEEE 802.11p protocols \cite{r4,r16} further
characterised the performance of different communication technologies
under realistic loads. Mir and Filali \cite{r16} conducted a useful
comparative study, though their results were simulation-based and the
specific case of lightweight embedded monitoring systems was not
addressed. Araniti et al.\ \cite{r4} analysed LTE for vehicular
networking, showing that cellular networks could handle the data
rates typical of real-time monitoring, providing a useful reference point for
understanding what is achievable with cloud transmission from a mobile
platform.

What is missing from all of this communication research is any
consideration of the application layer. Knowing that LTE can handle
10~Mbps tells you nothing about how to build a Flask server that
reliably pushes sensor readings to Firebase every 5 seconds. The
translation from network-level performance to practical embedded
system design was left as an exercise for the application developer.

\section{Navigation, Control Strategies, and IoT-Based Monitoring}

Yousif et al.\ \cite{r17} provided a comprehensive survey of visual
odometry and SLAM techniques, covering everything from monocular
systems to RGB-D approaches. The survey is a useful map of the field
but does not include implementations, so it is more useful as a
reference than as a practical guide.

Pacheco and Luo \cite{r10} ran a direct experimental comparison
between PID and Model Predictive Control on a physical mobile robot.
Their conclusion that PID performs well enough for structured
environments and is far easier to tune and maintain directly
informed the control strategy chosen in this work. Shaban and Nada
\cite{r9} reached similar conclusions comparing PID and PIP controllers,
reinforcing the case for PID as the practical default.

On the sensing side, Pushkaranand and Reddy \cite{r1} demonstrated an
IoT-based vehicular pollution monitoring system that could acquire
particulate matter and gas concentration data and push it to a web
interface in near real time. Ben-Aboud et al.\ \cite{r2} developed a
low-cost indoor air quality platform that logged CO$_2$, particulate
matter, and volatile organic compound data to a cloud backend. These
are direct predecessors of the monitoring subsystem described in
this thesis. Both systems classified air quality into categories
such as good, moderate, poor, and hazardous using CO$_2$ concentration
thresholds drawn from established indoor air quality standards.
Specifically, ASHRAE Standard 62.1 \cite{r46} and the WHO Indoor Air
Quality Guidelines \cite{r47} both define acceptable CO$_2$ levels for
occupied spaces and set thresholds above which ventilation is required.
Research by Sepp\"{a}nen et al.\ \cite{r48} showed that CO$_2$
concentrations above 1000\,ppm are associated with measurable health
effects in occupants, while work by Satish et al.\ \cite{r49}
demonstrated that even concentrations in the 1000--2500\,ppm range
impair cognitive decision-making performance. These findings provide
the scientific basis for the alert classification thresholds used in
the monitoring subsystem of this project.

Neither Pushkaranand et al.\ nor Ben-Aboud et al.\ included mobility.
Both systems were fixed installations. The data they collected was
representative of a single location, which is exactly the limitation
that motivated integrating sensing with a mobile platform here.

Pati and Kala \cite{r18} showed that vision-based robot following using
PID control was feasible on physical hardware, demonstrating that the
combination of visual feedback and proportional control is practical
and not just theoretical. Xu et al.\ \cite{r19} and Zhang et al.\
\cite{r20} added further experimental evidence on communication
performance and homography calibration for vision systems, respectively.

\section{Advances in Perception and Environmental Monitoring}

The introduction of deep learning into mobile robotics perception
changed the nature of the navigation problem. Siam et al.\ \cite{r12}
showed that semantic segmentation with convolutional networks could
achieve accuracy well beyond classical feature-based approaches for
path and lane detection. Fauadi et al.\ \cite{r21} surveyed
vision-based navigation more broadly, providing a useful taxonomy of
available techniques and their computational requirements.

Lozano Dom\'{i}nguez et al.\ \cite{r22} reviewed V2X, I2X, and P2X
communication technologies and highlighted the increasing need for
tight integration between sensing and communication systems, exactly
the problem that remains largely unaddressed in the literature at
a practical implementation level. Santana et al.\ \cite{r23} developed
vehicle-mounted air quality monitoring using IoT and MQTT, which is
probably the closest prior work to the framework described in this
thesis. Their platform collected particulate matter and gaseous
pollutant data from a moving vehicle and transmitted it to a cloud
backend. The key difference is that their system was designed around
a full-sized vehicle and focused on outdoor urban monitoring, not
the indoor structured-path scenario addressed here.

Dos Reis et al.\ \cite{r13} applied YOLO-based object detection to
autonomous navigation and achieved strong results, though the
computational cost was high for embedded deployment. Ben-Aboud et al.\
\cite{r2} provided further evidence that low-cost IoT platforms could
produce reliable air quality data, though mobility remained absent.

\section{Recent Developments in Integration and System-Level Approaches}

Papadeas et al.\ \cite{r24} surveyed real-time semantic segmentation
techniques and identified the trade-offs between accuracy and
computational cost that are particularly relevant for edge deployment.
Their finding that lighter models (such as nano-scale architectures)
can achieve near-comparable accuracy to much larger models at a
fraction of the compute cost directly supported the choice of
YOLOv8n-seg in this work.

Noordin et al.\ \cite{r11}, Singh et al.\ \cite{r26}, and Saidi
et al.\ \cite{r27} explored different control strategies for
differential-drive platforms, collectively confirming that PID remains
competitive despite its simplicity and that hybrid approaches offer
modest gains at substantially higher implementation complexity.

Shin et al.\ \cite{r28} and Petrov et al.\ \cite{r25} examined how
communication technology (including 5G and V2X protocols) performs
under real-world loads. Marroqu\'{i}n et al.\ \cite{r29} implemented
a complete embedded vision navigation system and provided a good
reference point for what is achievable on modern single-board computers
in terms of inference throughput and system latency.

More recent work from 2024 \cite{r30} has pushed further on perception
and path detection using DeepLab-based architectures, achieving very
high accuracy on controlled datasets. The trade-off, as with most
deep learning approaches, is model complexity and the need for a GPU
at inference time. For a Raspberry Pi~5 running at the edge, these
models are not practical without hardware accelerators.

\section{Research Gap}

Reading through this body of work, the fragmentation is the most
striking thing. There is solid work on navigation. There is solid work
on IoT monitoring. There is solid work on communication frameworks.
But no existing work, at least none that this review found,
combines all three into a single mobile platform that has been built,
deployed, and evaluated experimentally.

The specific gaps that this thesis addresses are:

\begin{itemize}
    \item Environmental monitoring systems in the literature are fixed.
          None of the reviewed systems mounts a sensor on a robot and
          drives a repeatable path through a space to collect spatially
          distributed data.
    \item Communication research operates at the network level and does
          not address the application-layer design decisions needed for
          real-time IoT deployment on embedded hardware.
    \item Navigation systems are evaluated without any environmental
          sensing payload. Adding a sensor subsystem changes the
          real-time constraints on the system and needs to be validated
          rather than assumed.
    \item No reviewed work separates the mobility subsystem from the
          sensing subsystem as independent modules that can be developed
          and maintained separately but deployed together.
    \item Multi-platform cloud storage, specifically writing simultaneously to a
          local database, a cloud backend, and an IoT analytics platform
          from a mobile embedded system, has not been demonstrated or
          validated in the reviewed literature.
\end{itemize}

These gaps directly motivate the design choices and experimental scope
of this work.

\section{Summary}

The evolution from early simulation-only navigation studies through to
embedded deep learning and IoT monitoring reflects both the maturation
of the individual fields and, more clearly, their continued isolation
from each other. Each area has developed impressive capabilities. What
has not kept pace is integration, and that is the problem this project
takes on.

The research gaps identified above guided every major design decision:
the choice of HTTP over MQTT, the separation of the Raspberry Pi~3
(monitoring) from the Raspberry Pi~5 (navigation), the use of a modular
software architecture, and the experimental protocol of 40 independent
trials across two operating modes. The literature survey is summarised
in Table~\ref{tab:litsurvey}.

% ── Literature Summary Table ─────────────────────────────────
\begingroup
\small
\setstretch{1.0}
\begin{longtable}{>{\centering\arraybackslash}p{0.7cm}
                  >{\centering\arraybackslash}p{0.7cm}
                  >{\RaggedRight\arraybackslash}p{2.8cm}
                  >{\RaggedRight\arraybackslash}p{2.8cm}
                  >{\RaggedRight\arraybackslash}p{2.8cm}
                  >{\RaggedRight\arraybackslash}p{2.7cm}}
\caption{Summary of Literature Survey.}
\label{tab:litsurvey}\\
\toprule
\textbf{Ref.} & \textbf{Year} & \textbf{Methodology} & \textbf{Achievements}
  & \textbf{Research Gap} & \textbf{Motivation} \\
\midrule
\endfirsthead
\toprule
\textbf{Ref.} & \textbf{Year} & \textbf{Methodology} & \textbf{Achievements}
  & \textbf{Research Gap} & \textbf{Motivation} \\
\midrule
\endhead
\bottomrule
\endfoot
%
\cite{r5}  & 1999 & Image-based visual servoing with EKF and nonlinear control
  & Established theoretical framework for vision-based robot navigation
  & Only simulation-based, no real implementation
  & Provides foundational base for mobile navigation \\ \hline
\cite{r6}  & 2002 & Homography-based visual servoing with Lyapunov stability
  & Developed stable control laws for robot motion
  & No experimental validation
  & Supports vision-based control concepts \\ \hline
\cite{r7}  & 2005 & Homography-based adaptive control
  & Achieved experimental validation of trajectory control
  & Requires structured environment
  & Provides practical control approach \\ \hline
\cite{r8}  & 2006 & Nonlinear adaptive control with homography tracking
  & Improved trajectory tracking accuracy
  & Needs predefined visual features
  & Useful for trajectory tracking understanding \\ \hline
\cite{r14} & 2007 & Shape-based visual servo control
  & Improved robustness of control system
  & Limited real-world benchmarking
  & Enhances control stability understanding \\ \hline
\cite{r3}  & 2008 & Secure communication protocols
  & Provided reliable communication architecture
  & No sensing integration
  & Basis for communication design \\ \hline
\cite{r15} & 2013 & MPC and Kalman filtering
  & Robust trajectory tracking under noise
  & High computational complexity
  & Advanced control reference \\ \hline
\cite{r4}  & 2013 & LTE-based communication analysis
  & Identified scalable communication solutions
  & No monitoring integration
  & Helps in communication understanding \\ \hline
\cite{r9}  & 2013 & PID vs PIP control comparison
  & Experimental validation of control methods
  & Limited test conditions
  & Supports control selection \\ \hline
\cite{r16} & 2014 & LTE vs 802.11p simulation
  & Compared communication technologies
  & No real-world implementation
  & Helps in choosing communication method \\ \hline
\cite{r17} & 2015 & SLAM and visual odometry survey
  & Comprehensive navigation techniques overview
  & No experimental validation
  & Provides background knowledge \\ \hline
\cite{r10} & 2015 & PID vs MPC experimental comparison
  & Demonstrated performance trade-offs
  & MPC complexity
  & Supports simpler control methods \\ \hline
\cite{r1}  & 2015 & IoT-based environmental monitoring
  & Enabled real-time sensing
  & No mobility integration
  & Direct motivation for monitoring system \\ \hline
\cite{r18} & 2017 & Vision-based robot following using PID
  & Real robot implementation
  & Limited evaluation
  & Shows practical feasibility \\ \hline
\cite{r19} & 2017 & DSRC vs LTE field experiments
  & Evaluated communication performance
  & No sensing integration
  & Helps communication understanding \\ \hline
\cite{r20} & 2017 & Homography calibration
  & Improved vision system accuracy
  & No monitoring integration
  & Supports perception accuracy \\ \hline
\cite{r12} & 2018 & CNN-based segmentation
  & Improved perception accuracy
  & High computational cost
  & Supports perception systems \\ \hline
\cite{r21} & 2018 & Vision navigation review
  & Provided system overview
  & No experimental validation
  & Background reference \\ \hline
\cite{r22} & 2019 & V2X communication survey
  & Identified communication trends
  & No sensing integration
  & Highlights research gap \\ \hline
\cite{r23} & 2019 & IoT + MQTT mobile monitoring
  & Real-time environmental monitoring
  & Limited integration
  & Directly relevant \\ \hline
\cite{r13} & 2019 & YOLO-based object detection
  & Improved navigation perception
  & High computational requirements
  & Supports perception \\ \hline
\cite{r2}  & 2020 & IoT monitoring platform
  & Low-cost sensing system
  & No mobility integration
  & Motivation for system design \\ \hline
\cite{r24} & 2021 & Segmentation survey
  & Identified trade-offs in perception methods
  & No implementation
  & Helps gap identification \\ \hline
\cite{r25} & 2021 & Communication benchmarking
  & Evaluated communication limits
  & Limited real-world testing
  & Supports communication analysis \\ \hline
\cite{r11} & 2021 & PID vs fuzzy control
  & Improved control performance
  & No real implementation
  & Supports control selection \\ \hline
\cite{r26} & 2022 & ML-based navigation
  & Improved navigation accuracy
  & Limited integration
  & Motivates integration \\ \hline
\cite{r27} & 2022 & Fuzzy-PID control
  & Improved system stability
  & Complex tuning
  & Motivates simpler approach \\ \hline
\cite{r28} & 2023 & V2X communication evolution
  & Improved communication efficiency
  & No sensing integration
  & Highlights communication gap \\ \hline
\cite{r29} & 2023 & Embedded vision navigation
  & Real-time system implementation
  & Limited benchmarking
  & Supports implementation \\ \hline
\cite{r30} & 2024 & Advanced perception \& control methods
  & High accuracy and performance
  & Mostly simulation-based and complex
  & Motivates simple integrated system \\
\end{longtable}
\endgroup